\section{Introduction}
\label{sec:introduction}

Virtual Machines (VMs) are a fundamental component of modern cloud computing, allowing multiple operating systems to run concurrently on a single physical machine. However, the performance of VMs can be significantly impacted by the overhead of memory protection mechanisms, particularly when it comes to network heavy workflows.

The Virtual Input/Output Memory Management Unit (vIOMMU) is a critical component of the virtualization stack that provides memory protection for I/O devices. It allows the hypervisor to control access to physical memory by VMs, ensuring that each VM can only access its own memory and preventing unauthorized access to other VMs' memory. However, the vIOMMU can introduce significant performance overheads, particularly in network-heavy workloads where frequent memory accesses are required. The overhead of the vIOMMU can be attributed to several factors, including the cost of page table walks, the cost of handling page faults, and the cost of managing the IOTLB (I/O Translation Lookaside Buffer) entries. The IOTLB is a cache that stores recently accessed memory translations for I/O devices, and it can become a bottleneck when there are frequent memory accesses that result in IOTLB misses

Traditionally, cloud providers had two options for managing the vIOMMU. They could either disable the vIOMMU, which would eliminate the overhead but also remove the memory protection for I/O devices allowing I/O devices to DMA to unprotected regions. They could also enable the vIOMMU and accept the performance overhead, such overheads are large as we will see. However, recent research has focused on developing optimized page unmapping strategies to mitigate the performance overhead of the vIOMMU while still providing memory protection. One such strategy is vFree ~\cite{vFree}, which aims to reduce the cost of IOTLB misses by optimizing the way pages are unmapped from the vIOMMU. 


